% !TeX root = ../main.tex
% Add the above to each chapter to make compiling the PDF easier in some editors.

\chapter{Introduction}\label{chapter:introduction}

This thesis develops and evaluates a pipeline that converts published histopathology literature into a structured knowledge base that is source-faithful and auditable, where every extracted rule, table, and figure can be traced back to its source in the paper. This chapter introduces and motivates that system; the following chapters build and measure it. 

Section~\ref{section:intro-motivation} explains the motivation behind the thesis; Section~\ref{section:intro-problem-statement} defines the problem, the gaps the thesis addresses, and the scope of the contribution. Section~\ref{section:intro-research-questions} presents the five research questions, and Section~\ref{section:intro-contributions} summarizes the contributions that answer them. Section~\ref{section:intro-approach-overview} gives a high-level overview of the entire pipeline. Section~\ref{section:intro-summary-of-findings} previews the main empirical findings, including the negative results of agreement-based cascading and their careful interpretation. Finally, Section~\ref{section:intro-thesis-structure} concludes the introduction by mapping the chapters that follow. 

\section{Motivation}\label{section:intro-motivation}

Biomedical knowledge is produced faster than any individual can read it, and much of this knowledge remains in text, tables and figures of full-text articles, rather than in a structured, machine-readable form. Open repositories such as PubMed Central make large collections of the existing literature available in text-form~\parencite{nlm_pmc_open_access}, making large-scale knowledge-extraction feasible in principle. However, in practice, diagnostically relevant facts, such as biomarker expressions of a tissue, the morphological features that support a diagnosis, or a finding is associated with better or worse prognosis, remain scattered in the text in the papers. 

Large language models (LLMs) make the automatic extraction of these findings from biomedical and clinical text more realistic~\parencite{}: a capable LLM can return typed claims with subject and outcome entities, as well as relations between these claims, more flexibly than traditional hand-built systems. This capability of LLMs is what makes it feasible to undertake the task of building a clinical knowledge-base over the existing histopathology literature. 

The same models, however, impose risks that a clinical knowledge-base cannot tolerate: they might hallucinate claims that seem to be correct at first glance, but are actually unsupported; omit important evidence, disregard nuance, or paraphrase claims in a way that makes the verification of these claims difficult~\parencite{}. Other work on this topic argues that the knowledge that an LLM extracts should be attributed, meaning that each claim should be tied to the evidence that supports it, so that an expert can inspect the evidence later~\parencite{}. 

This thesis therefore starts from an auditability requirement: the goal of the thesis is to extract claims from histopathology papers in a way that keeps each said claim traceable to its evidence. The intended output is a structured clinical knowledge base, the claims in which can be audited by an expert.
\section{Problem statement}\label{section:intro-problem-statement}

The components to build an auditable clinical knowledge base already exist, but mostly only in isolation. Layout-aware tools can reconstruct the logical structure of scientific PDFs~\parencite{}; dedicated detectors can detect tables and figures in PDFs, which often contain a lot of evidence in a histopathology paper~\parencite{}; biomedical NLI entity-normalization tools can map entities in different surface forms to shared biomedical concepts, so that findings across papers can be compared~\parencite{}; LLMs can perform structured extraction, and existing literature on the topic studies how to make the process financially more feasible, through cascading and routing mechanisms across multiple models~\parencite{}. Tracking where a piece of data came from to preserve provenance is also well-established~\parencite{}. 

The main problem is that these components are rarely combined into a single pipeline under a demanding cost-quality constraint; each component has its own benchmarks and its own definition of success. The objective of this thesis is more demanding than the evaluation of its components; the main goal is to build a pipeline that preserves complete provenance for all of its outputs, making each extracted rule traceable back to its source PDF and paragraph, while remaining economically feasible. 

The open question of this thesis is to address the question that asks what is required to bring together existing document-processing, biomedical normalization, LLM-extraction, NLI grounding, and model-cascading components, and see where off-the-shelf tools already suffice, where the calibrated agreement and grounding improve the results, and where they are worth their cost.

The scope of this thesis is bounded; the system targets open-access histopathology papers in English language, and is designed to emit a set of candidate clinical rules that will be reviewed and validated by experts. The confidence scores it emits should not be treated as clinical probabilities. 
\section{Research questions}\label{section:intro-research-questions}
\section{Contributions}\label{section:intro-contributions}
\section{Approach overview}\label{section:intro-approach-overview}

The system consists of two pipelines, connected by a PostgreSQL database, described in full in Chapter~\ref{chapter:Methodology}. The first is a document-extraction pipeline, which converts each PDF into hierarchical text, cropped figures and tables, as well as the page number, page coordinates, and section path each of these artifacts comes from. The document-extraction stage builds on a layout-aware extractor~\parencite{}, unions the results of the layout-aware extractor with that of the hybrid table detector~\parencite{}, and adds footnote expansion and two-pass invisible-text filtering. The provenance pointers of each element are stored in the database, so that the knowledge-extraction stages can use the source information for every piece of text. 

The second pipeline performs knowledge extraction. It turns the text emitted by the document-extraction pipeline into atomic, grounded, and normalized rules through a sequence of stages. The MAP stage produces typed claims, the grounding stage filters out claims that are not entailed by the evidence they cite, and entity normalization maps entities to shared biomedical concepts~\parencite{}. Afterwards, the grouping, canonicalization, relate, and resolve stages turn the surviving findings into ranked, canonicalized rule candidates, and emits the relations between them. Only the MAP stage is stochastic due to its use of language models; once the MAP output is fixed, all later stages are deterministic, and the pipeline outputs are reproducible. 

The cost-quality trade-off occurs in the MAP stage. Instead of sending each text chunk to the strongest model, an agreement-based cascading is used: cheaper voters from different providers summarize each chunk into the same typed schema, voter agreement is measured, and only the more difficult text chunks are escalated to the stronger voter tiers~\parencite{}. The same frozen biomedical NLI model~\parencite{} is used for grounding filtering and the relation classification. 
\section{Summary of findings}\label{section:intro-summary-of-findings}
\section{Thesis structure}\label{section:intro-thesis-structure}

The remainder of this thesis is structured as follows. 

Chapter~\ref{chapter:Related_Work} reviews existing research on scientific document extraction, biomedical information extraction and entity normalization, LLM-based structured extraction, natural language inference, provenance, cost-aware model cascading, and the evaluation of LLM-based extraction systems, then explains what the thesis adopts from prior work and where it diverges.

Chapter~\ref{chapter:Methodology} describes the system and the evaluation design. It presents the two-pipeline architecture, the corpus and how it was assembled, the document-extraction pipeline, the knowledge-extraction pipeline and its typed MAP-schema, the database backbone that connects the two pipelines, the agreement-based cascade, and the evaluation methodology that defines how the five research questions are measured. 

Chapter~\ref{chapter:Results} reports the empirical findings in the order of the research questions, from document extraction to held-out generalization.

Chapter~\ref{chapter:Discussion} interprets the results without introducing new measurements and discusses the main threats to validity. 

Chapter~\ref{chapter:Conclusion_and_Future_Work} summarizes the contributions and outlines future work, including human-annotated gold labels, generalization tests beyond the current corpus, and relation classification on real pipeline outputs, rather than on a synthetic dataset. 
